\chapter*{Abstract}
\addcontentsline{toc}{chapter}{Abstract}
\markboth{Abstract}{}

Small offices keep their records in a single, ever-growing spreadsheet. A coaching institute's fee
register repeats every teacher's phone number on hundreds of rows, spells the same student three
different ways, and needs an hour of filtering to answer ``who has not paid for October?''. These are
the update, insertion and deletion anomalies that normalization theory was invented to remove, but the
people who suffer from them have never heard the terms and the tools they use do not help: database
front-ends expect the user to design the tables, and chat-style assistants answer one question
without leaving a correct database behind.

QLens is a local web application that takes such a sheet and produces a normalized PostgreSQL database
from the data itself. The file is copied unchanged into a raw table and a typed staging table, with
column types inferred by SQL. Candidate keys and functional dependencies are then discovered inside the
database by PL/pgSQL functions (\code{qlens.check\_fd}, \code{qlens.is\_unique}) that group the staged
rows with dynamic SQL; a dependency that holds on all but a few rows is reported as a data
inconsistency together with the rows that break it. Near-duplicate spellings are found with
\code{pg\_trgm} trigram similarity over a GIN index, and every repair is logged. From the accepted exact
dependencies QLens computes a minimal cover and runs Bernstein's 3NF synthesis, extended with a key
relation so that the decomposition is always lossless, and turns the result into tables with surrogate
keys, \code{UNIQUE NULLS NOT DISTINCT} natural keys and foreign keys. The design is applied in a single
transaction and accepted only if a view that joins the new tables back together returns exactly the
staged rows, compared as multisets with \code{EXCEPT ALL} in both directions.

QLens keeps its own metadata in a normalized \code{qlens} schema of 22 tables with CHECK constraints,
generated columns, triggers and views that use window functions. Each workspace has its own schemas and
its own low-privilege PostgreSQL login role, so user queries are isolated by privileges rather than by
convention. On the bundled 3{,}000-row fee register the analysis finds the two composite candidate keys,
six core dependencies and the 13 planted faulty rows, and proposes the textbook four-table design
(\tbl{students}, \tbl{batches}, \tbl{teachers}, \tbl{fee\_register}); applying it rebuilds all 3{,}000
rows exactly. The algorithms are verified by unit tests on textbook examples, by 900 randomized synthesis
cases checked against a chase-based lossless-join oracle and a brute-force 3NF test, and by end-to-end
tests that apply the generated SQL to randomly generated sheets (302 backend and 63 frontend tests
pass). Applied designs are queried in a sandboxed, read-only SQL editor; a tool-using assistant is
in progress, and every capability is marked with its build status.

\medskip
\noindent\textbf{Keywords:} functional dependency discovery, third normal form, Bernstein synthesis,
lossless-join decomposition, PostgreSQL, PL/pgSQL, role-based isolation, data cleaning, \code{pg\_trgm}.
